% unidades/27-condicional-2.tex
\chapter{🔀 条件② — Condicional: たら・なら}
\unidadheader{27}{条件②}{Condicional たら・なら}{Hablar de situaciones hipotéticas y dar consejos}

\begin{tcolorbox}[vocabbox, title={📌 Vocabulario Nuevo}]
\begin{tabularx}{\linewidth}{llX}
\toprule
\textbf{Japonés} & \textbf{Español} & \textbf{Tipo} \\
\midrule
\vocabitem{もし}{—}{si (hipotético)}{conj.}
\vocabitem{〜場合}{ばあい}{en caso de ~ / si ~}{sust.}
\vocabitem{〜時}{とき}{cuando ~ / en el momento de ~}{sust.}
\vocabitem{暇}{ひま}{tiempo libre / sin plan}{adj-な/sust.}
\vocabitem{急に}{きゅうに}{de repente / súbitamente}{adv.}
\vocabitem{相談する}{そうだんする}{consultar / pedir consejo}{v. irr.}
\vocabitem{選ぶ}{えらぶ}{elegir / seleccionar}{v. gr.1}
\vocabitem{迷う}{まよう}{dudar / perderse}{v. gr.1}
\vocabitem{後悔する}{こうかいする}{arrepentirse}{v. irr.}
\vocabitem{チャンス}{—}{oportunidad}{sust.}
\bottomrule
\end{tabularx}
\end{tcolorbox}

\subsection*{🗣️ Frases Clave}

\frase{もし\ruby{宝}{たから}くじに\ruby{当}{あ}たったら、\ruby{何}{なに}をしますか？}
  {Si ganaras la lotería, ¿qué harías?}
\frase{\ruby{日本}{にほん}に\ruby{行}{い}くなら、\ruby{京都}{きょうと}がおすすめです。}
  {Si vas a Japón, te recomiendo Kioto.}
\frase{\ruby{困}{こま}ったら、\ruby{相談}{そうだん}してください。}
  {Si tienes problemas, consúltame.}

\subsection*{⚙️ Gramática}

\patron{Condición completada: たら}
  {V-た + ら (cuando/si se completa la acción)}
  {Más versátil: acciones, estados, pasado hipotético.
  \ej{\ruby{家}{いえ}に\ruby{帰}{かえ}ったら、\ruby{電話}{でんわ}してください。}
    {Cuando llegues a casa, llámame.}
  \ej{もし100\ruby{万円}{まんえん}あったら、\ruby{旅行}{りょこう}します。}
    {Si tuviera 1 millón de yen, viajaría.}}

\patron{Condición contextual: なら}
  {普通形 + なら (si es el caso que ~, en ese contexto)}
  {なら = recomendación basada en lo que el otro mencionó.
  \ej{\ruby{日本語}{にほんご}を\ruby{勉強}{べんきょう}するなら、アニメがいいですよ。}
    {Si vas a estudiar japonés, el anime es bueno.}
  \ej{あの\ruby{映画}{えいが}を\ruby{見}{み}るなら、\ruby{早}{はや}く\ruby{行}{い}ったほうがいい。}
    {Si vas a ver esa película, mejor ve pronto.}}

\begin{tcolorbox}[ejerciciobox, title={📝 Ejercicios Rápidos}]
\begin{enumerate}
  \item Usa たら para decir qué harás cuando termines de estudiar.
  \item Da un consejo con なら a alguien que quiere aprender a cocinar.
  \item ¿Cuál es la diferencia principal entre たら y なら?
\end{enumerate}
\vspace{0.4em}
\textbf{Respuestas:}
1) \ruby{勉強}{べんきょう}が\ruby{終}{お}わったら、(personal)\quad
2) \ruby{料理}{りょうり}を\ruby{勉強}{べんきょう}するなら、(recomendación personal)\quad
3) たら = acción/situación que se completa primero; なら = recomendación basada en contexto dado por el interlocutor.
\end{tcolorbox}

\begin{tcolorbox}[kanjibox, title={🀄 Kanjis de la Unidad}]
\begin{tabularx}{\linewidth}{cllXX}
\toprule
\textbf{Kanji} & \textbf{On} & \textbf{Kun} & \textbf{Significado} & \textbf{Vocab} \\
\midrule
\kanjirow{場}{じょう}{ば}{lugar / situación}{\ruby{場合}{ばあい} / \ruby{場所}{ばしょ}}
\kanjirow{急}{きゅう}{いそ・にわか}{urgente / de repente}{\ruby{急}{きゅう}に / \ruby{急行}{きゅうこう}}
\kanjirow{選}{せん}{えら}{elegir}{\ruby{選}{えら}ぶ / \ruby{選択}{せんたく}}
\kanjirow{迷}{めい}{まよ}{perderse / dudar}{\ruby{迷}{まよ}う / \ruby{迷路}{めいろ}}
\bottomrule
\end{tabularx}
\end{tcolorbox}
